%! Vectors — Unit 8, Lesson 8.6 — Homework Key
\documentclass[10pt]{article}
\usepackage{precalculus-article}
\usepackage{precalculus-key}
\newcommand{\TallMath}[1]{$\displaystyle #1\rule[-1.4em]{0pt}{3.2em}$}

\begin{document}

\pageheader{Unit 8, Lesson 8.6}{Homework Key}
\namedateperiod

\vspace{0.10in}

\begin{notesbox}{Vector Components and Magnitude}

\textbf{1.}\quad Find the component form and magnitude of the vector from
$P_1 = (-2, 1)$ to $P_2 = (4, 9)$.

\vspace{4pt}
$\vec{v} = \langle$ \ans{6}$,$ \ans{8}$\rangle$ \qquad
$\|\vec{v}\| = \ans{$\sqrt{6^2+8^2}=\sqrt{100}=10$}$

\begin{teachernote}
Component form: $\langle 4-(-2),\;9-1\rangle = \langle 6,8\rangle$.
Magnitude: $\sqrt{36+64}=10$.
\end{teachernote}

\vspace{14pt}
\textbf{2.}\quad A force vector has magnitude $r = 13$ and direction angle
$\theta = 67.4^\circ$. Given $\cos 67.4^\circ \approx 0.385$ and
$\sin 67.4^\circ \approx 0.923$, find the component form of the vector.
Round to one decimal place.

\vspace{4pt}
$a = r\cos\theta \approx \ans{$13(0.385)\approx 5.0$}$ \qquad
$b = r\sin\theta \approx \ans{$13(0.923)\approx 12.0$}$

\vspace{4pt}
$\vec{v} \approx \langle$ \ans{5.0}$,$ \ans{12.0}$\rangle$

\end{notesbox}

\vspace{0.08in}

\begin{notesbox}{Scalar Multiplication and Vector Addition}

\textbf{3.}\quad Let $\vec{u} = \langle 2, -3\rangle$ and $\vec{w} = \langle -4, 1\rangle$.
Compute each of the following.

\vspace{4pt}
(a)\quad $\vec{u} + \vec{w} = \ans{$\langle 2+(-4),\;-3+1\rangle = \langle -2,-2\rangle$}$

\vspace{8pt}
(b)\quad $3\vec{u} = \ans{$\langle 6,-9\rangle$}$

\vspace{8pt}
(c)\quad $2\vec{u} - \vec{w} = \langle 2(2)-(-4),\; 2(-3)-1 \rangle = \ans{$\langle 8,-7\rangle$}$

\vspace{8pt}
(d)\quad Express $\vec{u}$ in terms of $\hat{\imath}$ and $\hat{\jmath}$:
$\vec{u} = \ans{$2\hat{\imath} - 3\hat{\jmath}$}$

\end{notesbox}

\vspace{0.08in}

\begin{notesbox}{Unit Vectors and Dot Product}

\textbf{4.}\quad Find the unit vector in the direction of $\vec{v} = \langle -6, 8\rangle$.

\vspace{4pt}
$\|\vec{v}\| = \ans{10}$ \qquad
$\hat{u} = \left\langle \dfrac{-6}{\ans{10}},\; \dfrac{8}{\ans{10}}\right\rangle
= \langle$ \ans{$-\tfrac{3}{5}$}$,$ \ans{$\tfrac{4}{5}$}$\rangle$

\vspace{14pt}
\textbf{5.}\quad Let $\vec{p} = \langle 3, -4\rangle$ and $\vec{q} = \langle 4, 3\rangle$.

\vspace{4pt}
(a)\quad Compute $\vec{p} \cdot \vec{q}$: \quad
$(3)(4) + (-4)(3) = \ans{$12-12=0$}$

\vspace{8pt}
(b)\quad Are $\vec{p}$ and $\vec{q}$ perpendicular?
Circle one: \quad \textbf{Yes} \quad \textbf{No}

\vspace{4pt}
Justify: \ansline{Yes --- the dot product equals 0, which means the vectors are perpendicular.}

\end{notesbox}

\vspace{0.08in}

\begin{notesbox}{Angle Between Vectors and Vector Addition Triangle}

\textbf{6.}\quad Two vectors have magnitudes $\|\vec{u}\| = 5$ and $\|\vec{w}\| = 7$,
and the angle between them is $\theta = 60^\circ$.

\vspace{4pt}
(a)\quad Use the Law of Cosines to find the magnitude of $\vec{r} = \vec{u} + \vec{w}$.
(Recall: $\|\vec{r}\|^2 = \|\vec{u}\|^2 + \|\vec{w}\|^2 - 2\|\vec{u}\|\,\|\vec{w}\|\cos\theta$
where $\theta$ is the angle between $\vec{u}$ and $\vec{w}$.)

\vspace{4pt}
\ans{Angle in triangle $= 180^\circ - 60^\circ = 120^\circ$.\\[3pt]
$\|\vec{r}\|^2 = 25+49-2(5)(7)\cos(120^\circ) = 74+35 = 109$.\\[3pt]
$\|\vec{r}\| = \sqrt{109} \approx 10.4$.}

\begin{teachernote}
When the angle between vectors is $60^\circ$ as drawn from the same tail, the Law-of-Cosines triangle uses the supplement $120^\circ$. Result: $\|\vec{r}\|\approx10.44$.
\end{teachernote}

\vspace{8pt}
(b)\quad Given $\vec{a} = \langle 1, 2\rangle$ and $\vec{b} = \langle 2, -1\rangle$,
use the dot product formula $\theta = \cos^{-1}\!\left(\dfrac{\vec{a}\cdot\vec{b}}{\|\vec{a}\|\,\|\vec{b}\|}\right)$
to find the angle between $\vec{a}$ and $\vec{b}$. Show all steps.

\vspace{4pt}
\ans{$\vec{a}\cdot\vec{b}=(1)(2)+(2)(-1)=2-2=0$.\\[3pt]
$\|\vec{a}\|=\sqrt{5}$,\quad $\|\vec{b}\|=\sqrt{5}$.\\[3pt]
$\cos\theta=\dfrac{0}{\sqrt{5}\cdot\sqrt{5}}=0$,\quad so $\theta=\cos^{-1}(0)=90^\circ$.}

\end{notesbox}

\end{document}
